%% kapa: exam=1 round=19 year=2008 period=2 form=A subject=economics subject_ko=경제원론 seq=4 number=4
%% source: Q-Net 감정평가사 기출문제 (cst00309) artlSeq=1001962 file=extracted/1008/1001962/2008년 제19회 감정평가사 1차시험 기출문제/2교시.hwp
%% conversion: kordoc 4.24.51 -> tools/gen1.py (auto); review=proofread-sonnet; flags=none
\begin{question}{4}
솔로우(Solow) 경제성장모형에서 생산이 이루어지고 있을 때, 다음 설명 중 옳은 것을 모두 고르면?
\begin{framedtext}
ㄱ. 인구증가율이 낮아지면 1인당 국민소득은 장기적으로 낮아진다. \par
ㄴ. 저축률이 0.2에서 0.3으로 높아질 때 장기적으로 경제성장률이 높아진다. \par
ㄷ. 1인당 자본스톡(capital stock)이 균제상태(steady state)보다 낮은 수준에서 균제상태로 접근함에 따라 경제성장률이 둔화된다. \par
ㄹ. 자본소득을 모두 저축하고 노동소득을 모두 소비한다면, 이 경제의 균제상태는 황금률(golden rule)을 달성한다. \par
\end{framedtext}
\begin{choices}
  \item ㄱ, ㄴ
  \item ㄱ, ㄷ
  \item ㄷ, ㄹ
  \item ㄱ, ㄴ, ㄹ
  \item ㄴ, ㄷ, ㄹ
\end{choices}
\end{question}
\answer{3}
